% !PRINTED 256
\weberpage{256}
%<<CONT
mit alleiniger Ausnahme des Falles $n=1$, $p=2$. Zun\"achst haben wir in $G$ die Substitutionen
\[
\begin{pmatrix} 0, & -1 \\ 1, & 0 \end{pmatrix}\begin{pmatrix} 0, & 1 \\ -1, & \xi \end{pmatrix}\begin{pmatrix} 0, & 1 \\ -1, & 0 \end{pmatrix} = \begin{pmatrix} \xi, & 1 \\ -1, & 0 \end{pmatrix},
\]
also auch
\[
\begin{pmatrix} 0, & 1 \\ -1, & \xi \end{pmatrix}\begin{pmatrix} \xi, & 1 \\ -1, & 0 \end{pmatrix} = \begin{pmatrix} 1, & 0 \\ 2\xi, & 1 \end{pmatrix},
\]
und mithin
\[
\begin{pmatrix} 1, & 0 \\ 2\xi, & 1 \end{pmatrix}\begin{pmatrix} 0, & 1 \\ -1, & \xi \end{pmatrix} = \begin{pmatrix} 0, & 1 \\ -1, & 3\xi \end{pmatrix},
\]
\[
\begin{pmatrix} 1, & 0 \\ 2\xi, & 1 \end{pmatrix}\begin{pmatrix} 0, & 1 \\ -1, & 3\xi \end{pmatrix} = \begin{pmatrix} 0, & 1 \\ -1, & 5\xi \end{pmatrix} \ \text{u. s. f.,}
\]
\[
\begin{pmatrix} 0, & 1 \\ -1, & m\xi \end{pmatrix}
\]
f\"ur jede ungerade ganze Zahl $m$; wenn also $p$ ungerade ist, so k\"onnen wir $m=p$ annehmen und erhalten $A_0$. Ist aber $p=2$, so haben wir, wenn wir mit $\varrho$ eine noch zu bestimmende Zahl in $\mathfrak{C}$ bezeichnen, folgende Substitutionen in $G$:
\[
\begin{pmatrix} 0, & \varrho \\ -\varrho^{-1}, & 0 \end{pmatrix}\begin{pmatrix} 0, & 1 \\ -1, & \xi \end{pmatrix}\begin{pmatrix} 0, & -\varrho \\ \varrho^{-1}, & 0 \end{pmatrix} = \begin{pmatrix} \xi, & \varrho^2 \\ -\varrho^{-2}, & 0 \end{pmatrix}.
\]
Da hier nun jede Zahl ein Quadrat ist, so kann man $\varrho^2$ auch durch $\varrho$ ersetzen, und findet also, dass in $G$ die Substitution
\[
\begin{pmatrix} \xi, & \varrho \\ -\varrho^{-1}, & 0 \end{pmatrix}
\]
f\"ur jedes beliebige von Null verschiedene $\varrho$, und folglich auch die entgegengesetzte Substitution
\[
\begin{pmatrix} 0, & -\varrho \\ \varrho^{-1}, & \xi \end{pmatrix}
\]
vorkommen muss. Bedeuten also $\varrho,\sigma$ zwei von Null verschiedene Zahlen, so haben wir in $G$ auch
\weq{9}{\begin{pmatrix} 0, & \varrho^{-1}\sigma^{-1} \\ -\varrho\sigma, & \xi \end{pmatrix}\begin{pmatrix} \xi, & \sigma^{-1} \\ -\sigma, & 0 \end{pmatrix}\begin{pmatrix} 0, & \varrho \\ -\varrho^{-1}, & \xi \end{pmatrix} = \begin{pmatrix} 0, & 1 \\ -1, & \xi\varrho(1+\sigma+\sigma\varrho) \end{pmatrix}.}
